\documentclass[12pt]{article}
%% BUGS:
%% - The \begin{problem} ... \end{problem} is only used in exams??

\usepackage{url, graphicx}
\PassOptionsToPackage{normalem}{ulem} % magic
\usepackage{href-ul}

% page layout
\usepackage[letterpaper, textheight=9.5in, textwidth=5.5in]{geometry}
\usepackage{setspace}
\setstretch{1.08}
\pagestyle{plain}
\sloppy\sloppypar\raggedbottom\frenchspacing\thispagestyle{empty}

% problem set header
\newcommand{\psname}{Problem Set}
\newcommand{\startps}[1]{\section*{Computational Physics --- \psname~{#1}}}

% problem formatting
\usepackage{amsthm}
\theoremstyle{definition}
\newcommand{\problemname}{Problem}
\newtheorem{problem}{\problemname}
\newcommand{\probref}[1]{\problemname~\ref{#1}}
\theoremstyle{remark}
\newcommand{\notename}{Note}
\newtheorem{note}{\notename}[problem]
\newcommand{\noteref}[1]{\textit{\notename}~\ref{#1}}
\newcounter{subpart}[problem]
\newcommand{\subpart}{\par\addvspace{0.75ex}\refstepcounter{subpart}\textsl{(\alph{subpart})}~}

% words
\newcommand{\foreign}[1]{\textsl{#1}}
\newcommand{\vs}{\foreign{vs}}

% code
\newcommand{\code}[1]{\texttt{\detokenize{#1}}}
\usepackage{xcolor}
% Define custom colors for code highlighting
\definecolor{codegreen}{rgb}{0,0.6,0}
\definecolor{codegray}{rgb}{0.5,0.5,0.5}
\definecolor{codepurple}{rgb}{0.58,0,0.82}
\definecolor{backcolour}{rgb}{0.98,0.98,0.99}
\usepackage{listings}
\lstdefinestyle{pythonstyle}{
    backgroundcolor=\color{backcolour},   
    commentstyle=\color{codegreen},
    keywordstyle=\color{magenta},
    numberstyle=\tiny\color{codegray},
    stringstyle=\color{codepurple},
    tabsize=4
}
\lstset{
    basicstyle=\small\ttfamily,
    columns=fixed,
    breaklines=false,
    keepspaces=true,
    showspaces=false,
    showstringspaces=false,
    showtabs=false,
    showlines=*,
    style=pythonstyle,
    language=Python,
    literate={~}{{\url{~}}}1
}

% math
\usepackage{amsmath}
\newcommand{\dd}{\mathrm{d}}
\newcommand{\e}{\mathrm{e}}

% primary units
\newcommand{\rad}{\mathrm{rad}}
\newcommand{\kg}{\mathrm{kg}}
\newcommand{\m}{\mathrm{m}}
\newcommand{\s}{\mathrm{s}}

% secondary units
\renewcommand{\deg}{\mathrm{deg}}
\newcommand{\km}{\mathrm{km}}
\newcommand{\cm}{\mathrm{cm}}
\newcommand{\mm}{\mathrm{mm}}
\newcommand{\ft}{\mathrm{ft}}
\newcommand{\mi}{\mathrm{mi}}
\newcommand{\AU}{\mathrm{AU}}
\newcommand{\ns}{\mathrm{ns}}
\newcommand{\h}{\mathrm{h}}
\newcommand{\yr}{\mathrm{yr}}
\newcommand{\N}{\mathrm{N}}
\newcommand{\J}{\mathrm{J}}
\newcommand{\eV}{\mathrm{eV}}
\newcommand{\W}{\mathrm{W}}
\newcommand{\Pa}{\mathrm{Pa}}

% derived units
\newcommand{\mps}{\m\,\s^{-1}}
\newcommand{\mph}{\mi\,\h^{-1}}
\newcommand{\mpss}{\m\,\s^{-2}}

% random units
\newcommand{\au}{\mathrm{a.u.}}

% draft comments
\newcommand{\hogg}[1]{{\color{blue}\textbf{#1}}}

\begin{document}

\startps{6}
Due Friday 2026 October 16 at 18:00 on Brightspace.
Be sure to explicitly give credit where it is due, on all problems.
Also credit buddies, web pages, and LLMs (as per course policy).
Failure to give explicit credits will result in points deducted.

\begin{problem}
  In this problem we will build a model for the data we used in \psname~4, \problemname~1.
  The model is
  \begin{align}
    y_i = \beta + \alpha\,\exp(-\frac{1}{2}\,[x_i - \mu]^2) ~,
  \end{align}
  where $\beta,\alpha,\mu$ are free parameters.
  \subpart\label{sub:start}
  Write a function \code{predict(pars, xs)} that takes in as input \code{pars},
  which is a three-element array containing $\beta,\alpha,\mu$, and the \code{xs}
  that you loaded in from the pickle file.
  Plot the data that you plotted in \psname~4 in black squares, and overplot a model for the data
  in red exes or as a red line.
  Note that you will have to choose values for $\beta,\alpha,\mu$; choose them by hand so your plot looks ``okay.''
  \subpart
  Now make a function \code{loss(pars, xs, ys)} that takes in the same inputs, plus
  the $y$ values \code{ys}.
  Write this code such that it computes the total squared error---the sum of the squares of the difference between the $y$ values in \code{ys} and the ouptut of \code{predict}.
  This output should be just one number.
  Compute and report this output for the parameter values you used to make your plot in \subpartref{sub:start}.
  \subpart
  Now find and report the value of \code{pars} that minimizes \code{loss(pars, xs, ys)}.
  To do this you can use any optimizer you like.
  Look at \code{scipy.optimize} or write your own.
  Report your best-fit \code{pars} and re-make your plot from \subpartref{sub:start} but now with the optimized values.
  Does it look better?
\end{problem}

\begin{note}
  If you write your own optimizer, you might want to write \code{predict} and \code{loss} in Python Jax, so you get derivatives.
\end{note}

\begin{note}
  Your plots should be properly labeled here, but they don't have to be pretty.
  If you want to make the plots pretty, instead of plotting the prediction in a set of red exes or jagged line,
  plot it as a continuous function of $x$ by making a fine grid called \code{xplot},
  such that your plotting command looks like \code{plot(xplot, predict(pars, xplot), "r-")}.
\end{note}

\begin{problem}
  Newman, \textit{Computational Physics, 2ed}, Exercise~7.1.
\end{problem}

\begin{note}
  Make sure all the plots you make for this \problemname, and for all other \problemname s,
  meet Hogg's criteria for good plotting practice.
\end{note}

\begin{problem}
  Newman, \textit{Computational Physics, 2ed}, Exercise~7.2.
\end{problem}

\begin{problem}
  Newman, \textit{Computational Physics, 2ed}, Exercise~7.10.
\end{problem}

\end{document}
